%===============================================================================
% CHAPTER 2: LITERATURE REVIEW
%===============================================================================
\chapter{Literature Review and State of the Art}
\label{ch:literature}

\section{Traditional E-Voting Systems}
\label{sec:traditional-evoting}

\subsection{Direct Recording Electronic (DRE) Machines}
DRE machines replaced mechanical lever machines and punch cards in many jurisdictions
starting in the 1990s. While they improved usability and accessibility, numerous
studies \cite{kohno2004analysis, appel2009security} demonstrated critical vulnerabilities:
unverified software, lack of paper trails, and susceptibility to insider attacks.
The 2006 \textit{Princeton} study on the Diebold AccuVote-TS showed that malicious
firmware could alter votes undetectably \cite{felten2006security}.

\subsection{Internet Voting Systems}
Estonia's i-Voting system (operational since 2005) remains the most widely cited
internet voting deployment \cite{heiberg2011verifiability}. While it employs
cryptographic protocols for vote encryption and verification, researchers have
identified vulnerabilities in client-side malware resistance and procedural
security \cite{springall2014security}. Switzerland's Swiss Post system similarly
faced public scrutiny leading to a 2019 moratorium \cite{haenni2017swiss}.

\subsection{End-to-End Verifiable Voting}
Cryptographic protocols like \textit{Helios} \cite{adida2008helios}, \textit{Prêt à Voter}
\cite{ryan2009pret}, and \textit{Scantegrity} \cite{chaum2008scantegrity} introduced
mathematical verifiability: voters can confirm their vote was cast, recorded, and
tallied correctly without trusting any single party. However, these systems require
complex key ceremonies and have seen limited production adoption due to usability
barriers.

\section{Blockchain-Based Voting}
\label{sec:blockchain-voting}

\subsection{Smart Contract Platforms}
Early proposals used Ethereum smart contracts for vote recording \cite{mcmurray2018blockchain}.
The \textit{Follow My Vote} and \textit{VoteWatcher} projects demonstrated feasibility
but encountered practical barriers: high gas costs (often \$5--50 per vote during
congestion), 15--30 TPS throughput limiting election scale, and the requirement for
voters to manage cryptocurrency wallets and private keys.

Layer-2 solutions (Polygon, Arbitrum) reduce costs but introduce trust assumptions
about sequencer centralization. The fundamental tension remains: \textit{public
permissionless blockchains optimize for censorship resistance, not voting-specific
properties like receipt-freeness or coercion resistance}.

\subsection{Permissioned and Consortium Ledgers}
Hyperledger Fabric \cite{androulaki2018hyperledger} and Corda have been proposed
for organizational voting \cite{kumar2019blockchain}. These provide high throughput
and privacy but require a known validator set, shifting trust from a single
administrator to a consortium. For student elections or public organizations,
establishing a diverse validator set is often impractical.

\subsection{Hedera Hashgraph and HCS}
Hedera's Hashgraph consensus \cite{baird2016hashgraph} achieves asynchronous
Byzantine Fault Tolerance (aBFT) with fair ordering, providing mathematical
guarantees of consensus fairness. The Hedera Consensus Service (HCS) \cite{hedera2020hcs}
exposes this as a topic-based messaging layer: applications submit messages to
topics, and the network assigns consensus timestamps and ordering.

HCS advantages for voting:
\begin{itemize}
    \item \textbf{Fixed low fees:} \$0.0001 per message regardless of network load.
    \item \textbf{High throughput:} 10,000+ messages/second.
    \item \textbf{Fair ordering:} No front-running or transaction manipulation.
    \item \textbf{Immediate finality:} 3--5 seconds with mathematical certainty.
    \item \textbf{Public auditability:} Mirror nodes provide read-only access to history.
\end{itemize}

Recent work \cite{tran2023blockchain} has explored HCS for audit trails and voting,
demonstrating the architectural pattern this project implements.

\section{Hybrid Cloud-DLT Architectures}
\label{sec:hybrid-arch}

The \textit{Web2.5} paradigm \cite{chen2022web25} separates identity/access (Web2)
from consensus/state (Web3). AWS Serverless + Blockchain patterns have been
documented for supply chain \cite{aws2021blockchain} and identity \cite{aws2022identity},
but voting-specific architectures remain sparse.

This project's contribution is a complete, voting-optimized instantiation of the
hybrid pattern: identity via JWT/MongoDB, consensus via HCS, with careful attention
to the atomic coordination between the two layers (Section~\ref{sec:vote-flow}).

\section{Gap Analysis}
\label{sec:gap}

\begin{table}[H]
\centering
\caption{Comparison of Voting Approaches}
\label{tab:voting-comparison}
\begin{tabular}{@{}lccccc@{}}
\toprule
\textbf{Approach} & \textbf{Immutability} & \textbf{Cost/Vote} & \textbf{Throughput} & \textbf{Privacy} & \textbf{Complexity} \\
\midrule
Paper Ballots       & Physical only & \$2--10    & Manual    & Anonymous  & Low    \\
Centralized DB      & None          & Near zero  & 100k+ TPS & None       & Low    \\
Ethereum Smart C.   & \checkmark    & \$5--50    & 15--30    & Pseudonym. & High   \\
Hyperledger Fabric  & \checkmark    & Near zero  & 3,000+    & Config.    & High   \\
\textbf{This System (HCS)} & Prototype    & Not measured & Not measured & Pseudonym. & Medium \\
\bottomrule
\end{tabular}
\end{table}

The gap this project fills: \textbf{practical, low-cost, publicly verifiable voting
for real organizations} without requiring cryptocurrency expertise or consortium
governance.
